\section{Discussion}

The experimental results demonstrate that improvements in modern processor performance do not remove the need for efficient algorithms or memory-aware implementations. Instead, hardware improvements primarily change constant factors in execution time while leaving the underlying growth behavior determined by algorithmic complexity and memory access patterns.

The sorting experiment illustrates the continued importance of algorithmic complexity. For randomly ordered datasets of size \(n = 1{,}000\), Bubble Sort required approximately \(2.88\times10^{-2}\) seconds on average, compared to \(1.63\times10^{-4}\) seconds for \texttt{std::sort}. As the dataset size increased to \(n = 10{,}000\), Bubble Sort required approximately \(1.90\) seconds while \texttt{std::sort} required \(8.55\times10^{-4}\) seconds. For the largest dataset tested (\(n = 100{,}000\)), Bubble Sort required \(8.27\) seconds compared to \(1.73\times10^{-3}\) seconds for \texttt{std::sort}.

These measurements correspond to performance differences of approximately \(176\times\), \(2{,}223\times\), and \(4{,}792\times\), respectively. The widening gap between the algorithms closely reflects the theoretical difference between quadratic \(O(n^2)\) and \(O(n \log n)\) growth. Even on modern processors capable of executing billions of instructions per second, inefficient algorithms rapidly become impractical as input sizes increase.

The cache locality experiment highlights a second dimension of performance that is not captured by asymptotic complexity alone. Both traversal strategies perform the same number of arithmetic operations and therefore share identical theoretical complexity \(O(n^2)\) for traversing an \(N \times N\) matrix. However, the experiments reveal substantial runtime differences caused solely by memory access patterns.

For \(N = 1024\), column-major traversal was approximately \(2.49\times\) slower than row-major traversal. This slowdown increased to roughly \(4.20\times\) for \(N = 2048\) and \(5.78\times\) for \(N = 4096\). The increasing performance gap reflects the growing cost of cache misses as the working set size exceeds the capacity of faster cache levels. Because column-major traversal accesses memory with a stride proportional to the row length, consecutive accesses often fall into different cache lines, reducing spatial locality and increasing memory latency.

Taken together, these experiments illustrate two complementary principles of performance engineering. First, algorithmic complexity determines how runtime scales with input size. Second, the interaction between software and hardware architecture (particularly memory hierarchy behavior) can significantly affect performance even when algorithmic complexity remains unchanged. Modern hardware therefore does not eliminate the importance of efficient code; rather, it makes understanding both algorithm design and architectural behavior increasingly important.
